\section{Pasar apuntes a limpio en LaTeX}\label{sec:latex}

Los signos están dibujados con TikZ en \path{simbolos/glifos.tex}. No hace falta
ninguna fuente tipográfica ni imagen externa.

\subsection{Un signo dentro del texto}
\begin{Verbatim}[fontsize=\small,frame=single]
fricción \ks{K23} calor        % tamaño de apunte
fricción \kss{K23} calor       % más pequeño, para párrafos
\KalaSimbolo[20mm]{K23}        % cualquier tamaño
\end{Verbatim}
\KalaMostrar{fricción \ks{K23} calor \qquad fricción \kss{K23} calor \qquad \KalaSimbolo[20mm]{K23}}

\subsection{Una hoja de apunte}
El entorno \texttt{apunte} dibuja la hoja con su margen. Cada \verb|\linea| recibe
dos argumentos: lo que va en el margen y lo que va en el contenido. El título
es opcional.
\begin{Verbatim}[fontsize=\small,frame=single]
\begin{apunte}[Biología · célula]
\linea{}{mitocondria \ks{K27} célula eucariota}
\linea{\ks{K31}}{glucosa \ks{K24} ATP}
\end{apunte}
\lectura{La mitocondria es parte de la célula...}
\end{Verbatim}
\begin{apunte}[Biología · célula]
\linea{}{mitocondria \ks{K27} célula eucariota}
\linea{\ks{K31}}{glucosa \ks{K24} ATP}
\end{apunte}

\subsection{Añadir un signo}
Cada signo se registra una sola vez con \verb|\KalaRegistrar|: código, nombre,
familia, definición, número de trazos y dibujo. La cuadrícula es de
\(64\times64\) con el origen arriba a la izquierda; el dibujo debe quedar
aproximadamente entre 12 y 52.
\begin{Verbatim}[fontsize=\small,frame=single]
\KalaRegistrar{K36}{Nombre}{Conector}
  {Qué significa y qué no significa.}{2}{%
  \draw (12,32) -- (52,32);
  \draw (44,24) -- (52,32) -- (44,40);
}
\end{Verbatim}
Después se añade al catálogo (\path{docs/02-signos.tex}), al trazado
(\path{simbolos/trazado.tex}), a la hoja de consulta (\path{atlas.tex}) y al
cuaderno de práctica (\path{practica.tex}).

\subsection{Compilar}
Desde la carpeta del proyecto:
\begin{Verbatim}[fontsize=\small,frame=single]
make            # los tres PDF en compilado/
make manual     # solo este manual
make limpiar    # borra los archivos auxiliares
\end{Verbatim}
Se necesita pdf\LaTeX{} con TikZ y \texttt{latexmk}; no hace falta internet ni
\texttt{shell-escape}.
